% !PRINTED 287
\weberpage{287}
%<<CONT
in Übereinstimmung mit der Regel für die Zusammensetzung zweier Transformationen in § 28. Die Substitution (11) ist hiernach zu bezeichnen mit
\[
\begin{pmatrix} \partial, & c \\ b, & a \end{pmatrix}.
\]
Die durch alle Substitutionen dieser Form [nach (14)] gebildete Gruppe bezeichnen wir mit $\mathfrak{L}$ (Gruppe der linearen Substitutionen).

Die Funktionen (12) von $z'$ und also auch die Substitutionen (13) bleiben ungeändert, wenn die vier Zahlen $a$, $b$, $c$, $\partial$ mit einem und demselben durch $p$ nicht teilbaren Faktor multipliziert werden, und daraus ergibt sich nach (10) die Anzahl dieser Funktionen oder der Grad der Gruppe $\mathfrak{L}$
\weq{15}{p(p^2-1),}
den man auch leicht durch direkte Abzählung findet.

Werden in einer der Substitutionen (13) die vier Zahlen $a$, $b$, $c$, $\partial$ mit einem gemeinsamen, durch $p$ unteilbaren Faktor multipliziert, so wird die Determinante $\Delta$ mit dem \gesperrt{Quadrat} dieses Faktors multipliziert. Es bleibt daher nicht die Determinante $\Delta$, wohl aber ihr quadratischer Charakter, d.\,h. der Wert des Symbols
\[
\left(\frac{\Delta}{p}\right)
\]
durch diese Multiplikation erhalten.

Setzen wir zwei der Substitutionen (13) zusammen, so multiplizieren sich ihre Determinanten und hieraus folgt, daß \gesperrt{alle Substitutionen der Gruppe $\mathfrak{L}$, in denen $\Delta$ quadratischer Rest von $p$ ist, eine Gruppe unter sich bilden, die wir mit $\mathfrak{L}_0$ bezeichnen wollen.}

Die Gruppe $\mathfrak{L}_0$ ist ein (eigentlicher) Divisor der Gruppe $\mathfrak{L}$ vom Index 2.

Setzt man die Substitutionen von $\mathfrak{L}_0$ zusammen mit irgend einer Substitution $(z, \beta z) = \begin{pmatrix} 1, & 0 \\ 0, & \beta^{-1} \end{pmatrix}$, wo $\beta$ und also auch $\beta^{-1}$ ein quadratischer Nichtrest von $p$ ist, so erhält man die ganze Gruppe $\mathfrak{L}$.

Ist $\Delta$ quadratischer Rest von $p$, so kann man einen zu $a$, $b$, $c$, $\partial$ hinzuzufügenden gemeinsamen Faktor so wählen, daß $\Delta \equiv 1 \ (\mathrm{mod}\ p)$ wird, so daß wir die Gruppe $\mathfrak{L}_0$ auch darstellen können durch
%CONT>>
